\documentclass[10pt,a4paper]{article}
\usepackage[utf8]{inputenc}
\usepackage[T1]{fontenc}
\usepackage[french]{babel}
\usepackage[margin=1.3cm,top=1.1cm,bottom=1.1cm]{geometry}
\usepackage{titlesec}
\titlespacing*{\section}{0pt}{5pt}{1pt}
\titleformat*{\section}{\normalsize\bfseries}
\usepackage{amsmath,amssymb}
\usepackage{enumitem}
\setlist{nosep,leftmargin=1.2em}
\pagestyle{empty}
\setlength{\parindent}{0pt}
\setlength{\parskip}{2pt}
\setlength{\abovedisplayskip}{3pt}\setlength{\belowdisplayskip}{3pt}

\begin{document}

\begin{center}
{\Large\bfseries Modèle d'optimisation du dispatch des équipes en crise climatique}\\[2pt]
{\small Défi Enedis -- La Mêlée 2026 \quad·\quad moteur RSLS (règle de Smith + recherche locale, re-planifié à chaque événement)}
\end{center}

\section*{1. Ensembles et données}
\begin{itemize}
  \item $I$ : incidents (pannes) \emph{connus} à l'instant de calcul ; $K$ : équipes ($|K| = 63$, réparties sur 18 sites).
  \item $N_i$ : nombre de clients coupés par l'incident $i$ ; $s_i \in \{0,1\}$ : présence d'un client sensible.
  \item $r_i$ : heure de signalement de $i$ ; $p_i$ : durée de dépannage estimée ($2$, $4$ ou $8$ h).
  \item $\tau_{uv}$ : temps de trajet entre deux points $u,v$ (modèle calé sur 3\,718 trajets routiers réels).
  \item $\sigma_k$ : début du premier poste de l'équipe $k$ ; $W = 12$ h de travail, $R = 12$ h de repos.\\
        Fenêtres de travail de $k$ : $\mathcal{F}_k = \{[\sigma_k + n(W{+}R),\ \sigma_k + n(W{+}R) + W] : n \in \mathbb{N}\}$.
\end{itemize}

\section*{2. Variables de décision}
\begin{itemize}
  \item $x_{ik} \in \{0,1\}$ : vaut $1$ si l'équipe $k$ est affectée à l'incident $i$.
  \item $\pi_k = (i_1, i_2, \dots)$ : ordre dans lequel l'équipe $k$ traite ses incidents (une tournée par équipe).
  \item $S_i \ge 0$ : heure de départ de l'équipe vers $i$ ; \quad $C_i$ : heure de fin du dépannage de $i$.
\end{itemize}
Les heures ne sont pas choisies librement : un décodeur les calcule à partir de $(x,\pi)$ en plaçant chaque intervention
au plus tôt, dans une fenêtre de travail de l'équipe. Pour $i$ précédé de $j$ dans $\pi_k$ :
\[
S_i = \min\{\, t \ge \max(C_j,\ r_i) \;:\; [t,\ t + \tau_{ji} + p_i] \subseteq F \text{ pour un } F \in \mathcal{F}_k \,\},
\qquad C_i = S_i + \tau_{ji} + p_i .
\]
(pour la première intervention, $j$ désigne la position de l'équipe et $C_j$ l'heure à laquelle elle est libre).

\section*{3. Fonction objectif}
\[
\boxed{\;\min_{x,\ \pi}\; J = \sum_{i \in I} w_i\,(C_i - r_i) \;+\; \lambda_{\text{stab}} \cdot \Delta(\pi, \pi^{\text{préc}})\;}
\qquad\text{avec}\qquad w_i = N_i\,(1 + \beta\, s_i)\,\rho_i .
\]
\begin{itemize}
  \item $C_i - r_i$ : durée de coupure de l'incident $i$ (attente + trajet + dépannage).
  \item Mode officiel (indicateur du défi) : $\beta = 0$, $\rho_i = 1$, $\lambda_{\text{stab}} = 0$, soit
        $J = \sum_i N_i\,(C_i - r_i)$ = \textbf{clients $\times$ heures sans électricité}.
  \item Mode opérationnel : $\beta > 0$ donne plus de poids aux clients sensibles ; $\rho_i$ (priorité) et les autres
        règles viennent des refus du responsable ; $\Delta$ compte les équipes dont la prochaine mission change par
        rapport au plan précédent (stabilité du plan).
\end{itemize}

\section*{4. Contraintes}
\begin{enumerate}[label=(C\arabic*)]
  \item Chaque incident est traité par une seule équipe : $\sum_{k \in K} x_{ik} = 1 \quad \forall i \in I$.
  \item Une équipe fait une intervention à la fois, dans l'ordre $\pi_k$ (trajet entre deux interventions inclus).
  \item Pas d'anticipation : $S_i \ge r_i$ (on ne part pas vers une panne non signalée).
  \item Pas d'interruption : trajet + dépannage tiennent entièrement dans une fenêtre de travail de l'équipe.
  \item Les interventions déjà commencées sont figées ; les consignes du responsable (zone inaccessible jusqu'à $t$,
        équipe indisponible, interdiction ou affectation imposée, facteur de trajet) restreignent $(x, \pi)$.
\end{enumerate}

\section*{5. Contrôle de l'objectif de service (hors fonction objectif)}
L'objectif \og 90\,\% des incidents réparés en 48\,h \fg{} n'est \emph{pas} optimisé : il est vérifié sur le plan obtenu, avec $H_0$ le début de la crise :
\[
\text{part}_{48} = \frac{1}{|I|}\sum_{i \in I} \mathbf{1}\!\left[\, C_i \le H_0 + 48 \,\right] \;\ge\; 0{,}9 \ ?
\]
Si ce n'est pas le cas, le système cherche le plus petit nombre d'équipes de renfort $m$ (et leurs sites) tel que
$\text{part}_{48}(m) \ge 0{,}9$, et le propose au responsable.

\section*{6. Étape préalable : les horaires des équipes (fixés une fois, au début)}
Chaque équipe appartient à un groupe : A ($\sigma_k = 0$), B ($\sigma_k = 2$ h) ou C, la relève ($\sigma_k = 12$ h). Les effectifs
par site sont choisis sur 50 crises passées, puis évalués sur 50 autres :
\[
\min_{\sigma}\ \frac{1}{|\mathcal{C}|}\sum_{c \in \mathcal{C}} J_c(\sigma)
\qquad \text{s.c.}\qquad
\min_{t}\ \frac{1}{|K|}\,\big|\{k : t \in \textstyle\bigcup \mathcal{F}_k\}\big| \;\ge\; 20\,\%
\]
(au moins 20\,\% des équipes en service à tout instant ; repos de 12\,h, supérieur aux 11\,h légales).

\end{document}
